What is a robust representation of a risk measure?

Let us denote by \( \mathcal{M}_1 := \mathcal{M}_1(\Omega, \mathcal{F}) \) the set of all probability measures on \( (\Omega, \mathcal{F}) \).

Definition.
We say that a convex risk measure \( \rho \) has a robust representation if

\[
\rho(X) = \sup_{Q \in \mathcal{M}_1} \{ \mathbb{E}_Q[-X] - \alpha(Q) \} = \sup_{Q \in \mathcal{Q}_\rho} \{ \mathbb{E}_Q[-X] - \alpha(Q) \}
\]
with some penalty function \( \alpha : \mathcal{M}_1 \to \mathbb{R} \cup \{ \infty \} \). Here, the family

\[
\mathcal{Q}_\rho := \{ Q \in \mathcal{M}_1 \mid \alpha(Q) < \infty \}
\]

collects the measures \( Q \in \mathcal{M}_1 \) with finite penalty; in particular, for \(  Q \in \mathcal{Q}_\rho \), one requires that \( \mathbb{E}_Q[-X] \) 
is well-defined for all \( X \in \mathcal{X}. \) 

The measures in \(\mathcal{Q}_\rho\) can be interpreted as plausible probabilistic models that are taken more or less seriously as specified by the penalty function \(\alpha\).
The capital requirement \(\rho(X)\) is computed as the worst case of the penalized expected loss \(\mathbb{E}_Q[-X]\), taken over all models \(Q \in \mathcal{Q}_\rho\).


State the robust representation theorem for convex monetary risk measures for \( L^\infty \)
and \( L^p \). Discuss an economic interpretation in terms of Knightian uncertainty and
penalized losses.

Let \( \mathcal{M}_1(\mathbb{P}) ) \) denote all probability measures on which are abs. cont. with respect to \( \mathbb{P} \), so \( Q \ll \mathbb{P} \)
and the Radon-Nikodym derivative \( \frac{dQ}{d\mathbb{P}} \) exists.

Robust Representation of Risk Measures on \( L^\infty(\Omega, \mathcal{F}, \mathbb{P}) \)

Let \(\rho\) be a convex risk measure on \(\mathcal{X} = L^\infty(\Omega, \mathcal{F}, \mathbb{P})\). In this case, \(\rho\) admits a robust representation
\[
\rho(X) = \sup_{Q \in \mathcal{M}_1(\mathbb{P})} \left\{ \mathbb{E}_Q[-X] - \alpha(Q) \right\}
\]
with penalty function \(\alpha(Q) = \sup_{X \in A} \mathbb{E}_Q[-X]\), \(Q \in \mathcal{M}_1(\mathbb{P})\), \( A \) the acceptance set.
if and only if one of the equivalent conditions holds:

 \(\rho\) is continuous from above, i.e., \(\rho(X_n) \searrow \rho(X)\) whenever \(X_n \searrow X\) \(\mathbb{P}\)-a.s.
 \(\rho\) has the Fatou property:
\[
\rho(X) \leq \liminf_{n \to \infty} \rho(X_n)
\]
for any bounded sequence \((X_n)_{n \in \mathbb{N}}\) in \(L^\infty(\Omega, \mathcal{F}, \mathbb{P})\) which converges \(\mathbb{P}\)-a.s. to \(X\).

(Robust Representation of Risk Measures on \( L^p(\Omega, \mathcal{F}, \mathbb{P}) \)

A convex risk measure \(\rho : L^p(\Omega, \mathcal{F}, \mathbb{P}) \to \mathbb{R}\), \(1 \leq p < \infty\), admits a robust representation 
\[
\rho(X) = \max_{Q \in \mathcal{M}_1^q(\mathbb{P})} \{\mathbb{E}_Q[-X] - \alpha(Q)\}, \quad X \in L^p(\Omega, \mathcal{F}, \mathbb{P})
\]
with penalty function \(\alpha(Q) = \sup_{X \in A} \mathbb{E}_Q[-X]\), \(Q \in \mathcal{M}_1^q(\mathbb{P})\), with
- dual exponent \(q := p/(p-1) \in (1, \infty]\),
- probability measures \(\mathcal{M}_1^q := \left\{ Q \in \mathcal{M}_1(\mathbb{P}) \bigg| \frac{dQ}{d\mathbb{P}} \in L^q(\Omega, \mathcal{F}, \mathbb{P}) \right\}\).

What are the following objects: Young function, Luxemburg norm, Orlicz space, Orlicz heart?

The question goes on with:
How are the latter two space related to each other? State the \( \Delta_2 \) -
condition and explain its significance. What is the connection to \( L^p \) -spaces? Can
you state a robust representation theorem for monetary risk measures defined on
Orlicz hearts?

A Young function \(\Psi: [0, \infty) \to [0, \infty)\)
is left-continuous, non-decreasing convex such that \(\lim_{x \to 0} \Psi(x) = 0\), \(\lim_{x \to \infty} \Psi(x) = \infty\).

The Orlicz space
\(
L^\Psi := L^\Psi(\Omega, \mathcal{F}, \mathbb{P}) = \{X \in L^0 : \mathbb{E}[\Psi(c|X|)] < \infty \text{ for some } c > 0 \}
\)
is a Banach space if endowed with the Luxemburg norm
\(
\|X\|_\Psi := \inf\{a > 0 : \mathbb{E}[\Psi(|X|/a)] \leq 1\}.
\)

The Orlicz heart
\(H^\Psi := H^\Psi(\Omega, \mathcal{F}, \mathbb{P}) = \{X \in L^0 : \mathbb{E}[\Psi(c|X|)] < \infty \text{ for all } c > 0\}\)
is a linear subspace of \(L^\Psi\) that equals the trivial space \(\{0\}\), if \(\Psi\) attains the value \(\infty\). 
We will, thus, assume from now on that \(\Psi\) is a real-valued function. In this case, \(\Psi\) is continuous,
\(
L^\infty \subseteq H^\Psi \subseteq L^\Psi \subseteq L^1,
\)
and the Orlicz heart \(H^\Psi\) is the closure of \(L^\infty\) in \(L^\Psi\) with respect to the Luxemburg norm \(\|\cdot \|_\Psi\), thus also a Banach space.

The equality \(H^\Psi = L^\Psi\) holds, if and only if the \(\Delta_2\)-condition holds:
\(
\exists C, x_0 > 0 \, \forall x \geq x_0 : \Psi(2x) \leq C \Psi(x).
\)

An important special case in which the \(\Delta_2\)-condition is satisfied corresponds to Young functions \(\Psi(x) = x^p/p\) for \(p \in [1, \infty)\), 
leading to the classical situation of \(L^p\)-spaces, i.e.,
\(
L^p = H^\Psi = L^\Psi.
\)
In contrast, if we consider the exponential Young function \(\Psi(x) = e^x - 1\) which does not satisfy the \(\Delta_2\)-condition, 
the Orlicz heart and Orlicz space are not identical.

Suppose that a convex risk measure \(\rho : H^\Psi \to \mathbb{R}\) is defined on the Orlicz heart \(H^\Psi\), then it admits a robust representation
\(
\rho(X) = \max_{Q \in \mathcal{M}^{\Psi^\ast}(\mathbb{P})}\{\mathbb{E}_{\mathbb{Q}}[-X] - \alpha(\mathbb{Q})\}
\)
where \( X \in H^\Psi \)
in terms of the set
\(\mathcal{M}^{\Psi^\ast}(\mathbb{P}) := \{ \mathbb{Q} \in \mathcal{M}_1(\mathbb{P})| \frac{d\mathbb{Q}}{d\mathbb{P}} \in L^{\Psi^\ast}(\mathbb{P}) \}\)
where the conjugate of \( \Psi \)
\( \Psi^\ast(y) := \sup_{x \geq 0} \{ xy - \Psi(x) \},\ y\geq 0\)
is again a Young function.


Distribution-based convex risk measures possess a unique extension to \( L^1 \). Explain
this in detail. How is the finiteness of the risk measure on subspaces related to
continuity properties?

Distribution-Based Risk Measure
A monetary risk measure \(\rho\) on \(\mathcal{X}\) is called distribution-based (or law-invariant), if \( \rho(X) \) depends only on the distribution of \( X \) under \(\mathbb{P}\), i. e., \(\rho(X)=\rho(Y)\), whenever \( X \) and \( Y \) have the same distribution under \(\mathbb{P}\).

Theorem.
Let \(\rho:L^{\infty}\rightarrow\mathbb{R}\) a distribution-based convex risk measure. 
- Then there exists a unique extension 
\[\bar{\rho}:L^{1}\rightarrow\mathbb{R}\cup\{+\infty\}\]
   of \(\rho\) that is convex, monotone, cash-additive, and lower semicontinuous with respect to the \(L^{1}\)-norm.
- If \(\bar{\rho}\) is finite on some Orlicz heart \(H^{\omega}\) with finite Young function \(\Psi\), then it is continuous on \(H^{\omega}\) with respect to the corresponding Luxemburg norm.

Risk measures can thus uniquely be defined on either \(L^{\infty}\) or \(L^{1}\). Their continuity properties are related to their finiteness.


Compute the value at risk of a lognormal loss.


\(Y \sim N(\mu, \sigma^2)\)
Financial position \(X = \exp(Y)\), whose losses \(L = -X\) have a lognormal distribution
Note that both \(X\) and \(Y\) have continuous distributions.

Derivation via quantile transformation:
\[V@R_{\lambda}(X) = -q_X(\lambda) = \exp(q_Y(1-\lambda)) = \exp(V@R_{\lambda}(-Y)) = \exp(\mu + \Phi^{-1}(1-\lambda)\sigma)\]

Alternative derivation:
The Value at Risk of \(X\) at level \(\lambda\) is determined by:
\[\lambda = P[X + V@R_{\lambda}(X) < 0] = P[-\exp(Y) < -V@R_{\lambda}(X)] = P[Y > \ln(V@R_{\lambda}(X))] = 1 - \Phi\left(\frac{\ln(V@R_{\lambda}(X)) - \mu}{\sigma}\right).\]
This implies:
\[\Phi^{-1}(1-\lambda) = \frac{\ln(V@R_{\lambda}(X)) - \mu}{\sigma} \Leftrightarrow V@R_{\lambda}(X) = \exp(\mu + \Phi^{-1}(1-\lambda)\sigma)\]

What is the statement of Kusuoka’s theorem?

Let \(M_1((0,1])\) denote the set of all probability measures on \(((0, 1], \mathcal{B}((0, 1]))\), 
where \(\mathcal{B}((0, 1])\) is the Borel-\(\sigma\)-algebra on the interval \((0, 1]\).

Robust Representation of Distribution-Based Risk Measures:
Suppose that \(L^2(\Omega, \mathcal{F}, P)\) is separable. In this case, a convex risk measure \(\rho: L^{\infty}(\Omega, \mathcal{F}, P) \rightarrow \mathbb{R}\) with acceptance set 
\(A_{\rho}\), is distribution-based if and only if
\[\rho(X) = \sup_{\mu \in M_1((0,1])} \left( \int_{(0,1]} AV@R_{\lambda}(X) \mu(d\lambda) - \beta_{\min}(\mu) \right)\]
with penalty function \(\beta_{\min}(\mu) = \sup_{X \in A_{\rho}} \int_{(0,1]} AV@R_{\lambda}(X) \mu(d\lambda)\).



TODO 

Define the risk measures average value at risk. State key properties of this risk
measure. Prove that it dominates V@R at the same level. What is its robust
representation?
• State its OCE-representation?
(This is particularly important for Markowitz-type portfolio problems!)


Derive a formula for the AV@R of normal random variables. How does this compare
to V@R?

AV@R and variance: If \(X\) is normally distributed, then
\[AV@R_{\lambda}(X) = E_P[-X] + \frac{1}{\lambda} \varphi(\Phi^{-1}(\lambda)) \sigma_P(X),\]
where \(\varphi\) denotes the density of the standard normal distribution.

Derivation:
Use that \(V@R_{\lambda}(X) = E_P[X] - \Phi^{-1}(\lambda) \sigma_P(X)\) combined with the substitution \(x = \Phi^{-1}(\alpha)\) and \(\varphi'(x) = -x\varphi(x)\):

\[
AV@R_{\lambda}(X) &= \frac{1}{\lambda} \int_0^{\lambda} V@R_{\alpha}(X) d\alpha = E_P[-X] + \left( \frac{1}{\lambda} \int_0^{\lambda} -\Phi^{-1}(\alpha) d\alpha \right) \sigma_P(X)
\]
\[
= E_P[-X] + \left( \frac{1}{\lambda} \int_{-\infty}^{\Phi^{-1}(\lambda)} -x\varphi(x) dx \right) \sigma_P(X) \\
\]
\[
= E_P[-X] + \frac{1}{\lambda} \varphi(\Phi^{-1}(\lambda)) \sigma_P(X)
\]
